% ---------------------------------------------------------------------------
% Capítulo de Metodologia (rascunho)
% Base: docs/rsl1-protocolo.md (v0.22), docs/rsl2-protocolo.md (v0.18),
% docs/buscas/rsl1-definitivas-2026-11-10.md, docs/rsl1-calibracao.md,
% docs/requisitos-app.md e docs/referencias-verificadas.md.
% Situação descrita: 12/11/2026.
% ---------------------------------------------------------------------------

% As citações longas (por exemplo, KITCHENHAM; MADEYSKI; BUDGEN) não quebram
% entre linhas; o espaço de emergência evita que saiam da margem. O grupo
% limita o ajuste a este capítulo.
\begingroup
\setlength{\emergencystretch}{3em}

\chapter{Metodologia}
\label{cap:metodologia}

Este capítulo descreve como o trabalho é conduzido: a classificação da
pesquisa, as etapas do trabalho, o método das duas revisões sistemáticas da
literatura (RSLs), a forma de registro que torna as revisões verificáveis e o
levantamento de requisitos do buscador. As etapas ainda não executadas são
descritas como procedimento.

% ===========================================================================
\section{Classificação da pesquisa}
\label{sec:metodologia-classificacao}

% PENDENTE: classificar a pesquisa quanto à natureza, à abordagem, aos
% objetivos e aos procedimentos. Nenhuma referência de
% docs/referencias-verificadas.md trata de classificação de pesquisa.
% É preciso escolher uma obra de metodologia científica revisada por pares,
% conferi-la e registrá-la antes de escrever esta seção.

O trabalho reúne dois tipos de estudo. O primeiro são duas revisões
sistemáticas da literatura, que são estudos secundários e seguem diretrizes de
relato próprias para esse tipo de estudo
\cite{page2021prisma,kitchenham2023segress}. O segundo é o desenvolvimento de
um protótipo de buscador de inteligência de fontes abertas (OSINT) para órgãos
policiais e de inteligência, cujos requisitos partem de um levantamento com
especialista e da evidência reunida na primeira revisão.

% ===========================================================================
\section{Etapas do trabalho}
\label{sec:metodologia-etapas}

O trabalho tem quatro etapas, ligadas entre si:

\begin{enumerate}
	\item \textbf{Revisões sistemáticas da literatura.} A RSL 1 trata de
	ferramentas OSINT para órgãos policiais e de inteligência. Cada achado
	relevante dela vira um requisito candidato do buscador. A RSL 2 trata do
	desenvolvimento de software com modelos de linguagem de grande porte
	(LLMs) e agentes de IA. Cada achado relevante dela vira uma prática ou um
	requisito candidato para o processo de desenvolvimento do buscador.
	\item \textbf{Levantamento de requisitos com especialista.} Uma lista
	inicial de requisitos funcionais e não funcionais, recebida do orientador,
	é confrontada com a evidência da RSL 1 e com a legislação
	(Seção~\ref{sec:metodologia-requisitos}).
	\item \textbf{Desenvolvimento do protótipo.} O protótipo implementa o
	recorte dos requisitos definido na etapa anterior.
	\item \textbf{Avaliação.} Os critérios de avaliação do protótipo partem da
	QP5 da RSL 1, que trata de como as ferramentas são avaliadas e de como
	garantem a validade e a preservação da evidência coletada.
\end{enumerate}

% PENDENTE: descrever o processo de desenvolvimento do protótipo (ciclo,
% práticas, ferramentas e uso de assistentes de IA, ligado à RSL 2). O
% recorte do protótipo e a stack ainda não foram definidos
% (docs/requisitos-app.md, "Pendências"). Não há referência verificada sobre
% processo de desenvolvimento.

% PENDENTE: definir o método de avaliação do protótipo (critérios, métricas,
% participantes, procedimento). Depende dos resultados da QP5 da RSL 1 e de
% uma referência verificada sobre avaliação.

% ===========================================================================
\section{Método das revisões sistemáticas}
\label{sec:metodologia-rsl}

As duas revisões são revisões sistemáticas completas e seguem o mesmo método.
Esta seção apresenta o método comum e aponta as diferenças entre elas.

% ---------------------------------------------------------------------------
\subsection{Referenciais metodológicos}

As revisões seguem quatro referenciais:

\begin{itemize}
	\item \textbf{Relato:} a declaração PRISMA 2020 \cite{page2021prisma} e,
	para a estratégia de busca, a extensão PRISMA-S \cite{rethlefsen2021}.
	\item \textbf{Condução:} o documento de explicação e elaboração do PRISMA
	2020 \cite{page2021ee}, que descreve, item a item, as formas de conduzir
	cada etapa.
	\item \textbf{Relato em engenharia de software:} a diretriz SEGRESS,
	que estende os itens do PRISMA 2020 a estudos secundários da área,
	inclusive estudos de mapeamento e revisões qualitativas
	\cite{kitchenham2023segress}.
	\item \textbf{Busca complementar:} o \textit{snowballing}
	\cite{wohlin2014}.
\end{itemize}

% ---------------------------------------------------------------------------
\subsection{Objetivos e escopo}

O objetivo da RSL 1 é identificar e caracterizar, na literatura revisada por
pares, as ferramentas, técnicas e requisitos de sistemas de busca OSINT
voltados a órgãos policiais e de inteligência, incluindo o uso de inteligência
artificial e os desafios legais, éticos e de segurança associados.
\citeonline{ghioni2024} tratam das implicações de governança, ética, direito e
sociedade da OSINT com IA, e \citeonline{browne2024}, das aplicações de IA à
OSINT em geral. O trabalho de \citeonline{ghioni2024} é uma revisão de
literatura com análise bibliométrica, montada com busca por título no Google
Scholar e no Scopus, e não uma revisão no formato PRISMA. Nenhuma das duas
revisões restringe o escopo ao uso por órgãos policiais e de inteligência, que
é o recorte da RSL 1.

O objetivo da RSL 2 é identificar e sintetizar as evidências empíricas,
publicadas em veículos revisados por pares, sobre os efeitos do uso de
assistentes e agentes baseados em LLMs no desenvolvimento de software, nas
dimensões de produtividade, qualidade, segurança, custo e processo. Três
revisões têm recortes próprios: \citeonline{hou2024} analisaram 395 artigos,
de janeiro de 2017 a janeiro de 2024, sobre a aplicação de LLMs a tarefas de
engenharia de software; \citeonline{liu2026} reuniram 124 artigos sobre
agentes baseados em LLMs na engenharia de software; e \citeonline{mohamed2026}
revisaram 39 estudos, de janeiro de 2014 a dezembro de 2024, sobre o impacto
de assistentes baseados em LLMs na produtividade dos desenvolvedores. A RSL 2
junta as dimensões das suas questões de pesquisa num mesmo corpus e inclui
estudos publicados depois de 2024.
% PENDENTE: os números e os períodos de Hou et al., Liu et al. e Mohamed et
% al. foram conferidos só no resumo, na página do editor
% (referencias-verificadas.md, "Recortes das revisões e lacuna").

\textbf{Terminologia da RSL 2.} O trabalho chama de desenvolvimento
\textit{AI-First} o desenvolvimento de software em que assistentes e agentes
baseados em LLMs são o meio principal de produzir artefatos (código, testes,
documentação), sob supervisão humana. Essa é uma definição operacional deste
trabalho. Numa busca exploratória por título na Crossref, em 06/11/2026, o
termo apareceu em 12 registros, e nenhum artigo de periódico entre eles o usa
para o processo de desenvolvimento. Por isso o termo \textit{AI-First} não
entra na string de busca, que usa os termos das três revisões citadas acima:
LLMs \cite{hou2024}, agentes baseados em LLMs \cite{liu2026} e assistentes
baseados em LLMs \cite{mohamed2026}, além de IA generativa.

O escopo das duas revisões segue a estrutura PICO, que o PRISMA 2020 indica
para formular os objetivos de uma revisão \cite[item 4]{page2021ee}, acrescida
do contexto (C). A Tabela~\ref{tab:picoc} resume o escopo de cada revisão.

\begin{table}[htbp]
	\centering
	\caption{Escopo das revisões (PICOC)}
	\label{tab:picoc}
	\small
	\begin{tabular}{p{2.1cm}>{\raggedright\arraybackslash}p{6.1cm}>{\raggedright\arraybackslash}p{6.1cm}}
		\hline
		\textbf{Elemento} & \textbf{RSL 1} & \textbf{RSL 2} \\
		\hline
		População & Órgãos policiais, de investigação criminal e de
		inteligência; analistas e investigadores &
		Desenvolvimento de software e quem o faz: desenvolvedores, equipes e
		projetos (industriais, de código aberto ou acadêmicos) \\
		\hline
		Intervenção & Ferramentas, sistemas, plataformas, arquiteturas e
		métodos de coleta e análise OSINT, com ou sem IA &
		Uso de assistentes e agentes baseados em LLMs em atividades de
		desenvolvimento, em qualquer nível de autonomia \\
		\hline
		Comparação & Não se aplica &
		Desenvolvimento sem IA, ou com outra ferramenta, modelo ou nível de
		autonomia, quando o estudo fizer essa comparação \\
		\hline
		Resultados & Funcionalidades, arquiteturas, técnicas de IA,
		requisitos, desafios, ameaças e formas de avaliação &
		Produtividade, qualidade do software, segurança, custo e mudanças no
		processo e na governança \\
		\hline
		Contexto & Investigação criminal, segurança pública e atividade de
		inteligência &
		Indústria, código aberto e academia \\
		\hline
	\end{tabular}
	\fonte{Elaborada pelo autor, com base nos protocolos das revisões.}
\end{table}

% ---------------------------------------------------------------------------
\subsection{Questões de pesquisa}

As Tabelas~\ref{tab:qp-rsl1} e~\ref{tab:qp-rsl2} apresentam as questões de
pesquisa (QP) de cada revisão e a parte do trabalho que cada uma alimenta.

\begin{table}[htbp]
	\centering
	\caption{Questões de pesquisa da RSL 1}
	\label{tab:qp-rsl1}
	\small
	\begin{tabular}{p{0.9cm}p{9.2cm}>{\raggedright\arraybackslash}p{4.2cm}}
		\hline
		\textbf{ID} & \textbf{Questão} & \textbf{Uso no trabalho} \\
		\hline
		QP1 & Quais ferramentas, sistemas e arquiteturas de busca OSINT foram
		propostos ou avaliados para uso policial e de inteligência? &
		Estado da arte e referência para a arquitetura do buscador \\
		\hline
		QP2 & Quais técnicas de IA (incluindo modelos de linguagem) são
		usadas, e em quais fases do ciclo de inteligência? &
		Onde a IA agrega valor no buscador \\
		\hline
		QP3 & Quais requisitos e desafios legais, éticos, de privacidade e de
		soberania de dados são relatados? &
		Requisitos não funcionais \\
		\hline
		QP4 & Quais ameaças de segurança e riscos são identificados
		(inclusive contra a ferramenta e contra o investigador), e como são
		mitigados? &
		Modelo de ameaças \\
		\hline
		QP5 & Como as ferramentas são avaliadas, e como garantem a validade e
		a preservação da evidência coletada? &
		Critérios de avaliação e cadeia de custódia \\
		\hline
	\end{tabular}
	\fonte{Elaborada pelo autor, com base no protocolo da RSL 1.}
\end{table}

\begin{table}[htbp]
	\centering
	\caption{Questões de pesquisa da RSL 2}
	\label{tab:qp-rsl2}
	\small
	\begin{tabular}{p{0.9cm}p{9.2cm}>{\raggedright\arraybackslash}p{4.2cm}}
		\hline
		\textbf{ID} & \textbf{Questão} & \textbf{Uso no trabalho} \\
		\hline
		QP1 & Quais efeitos sobre a produtividade dos desenvolvedores são
		relatados, e com quais métricas? &
		Expectativa de ganho no desenvolvimento do buscador \\
		\hline
		QP2 & Quais efeitos sobre a qualidade do software (defeitos,
		manutenibilidade, testes, legibilidade) são relatados? &
		Práticas de qualidade no desenvolvimento \\
		\hline
		QP3 & Quais vulnerabilidades e riscos de segurança o código gerado e
		os agentes introduzem, e quais mitigações foram avaliadas? &
		Requisitos de segurança e revisão de código \\
		\hline
		QP4 & Quais custos (financeiros, computacionais, de \textit{tokens},
		de energia) são relatados, e como são medidos? &
		Custo de desenvolver com IA \\
		\hline
		QP5 & Como o processo de engenharia e a governança mudam (papéis,
		revisão humana, etapas do ciclo de vida, políticas de uso)? &
		Processo de desenvolvimento do buscador \\
		\hline
		QP6 & Quais ferramentas, modelos e níveis de autonomia (completar
		código, assistente conversacional, agente autônomo) foram estudados,
		e em quais contextos? &
		Mapeamento do corpus; contextualiza as QP1 a QP5 \\
		\hline
	\end{tabular}
	\fonte{Elaborada pelo autor, com base no protocolo da RSL 2.}
\end{table}

% ---------------------------------------------------------------------------
\subsection{Fontes de busca}

\textbf{Bases internacionais.} As duas revisões buscam em quatro bases:
Scopus, IEEE Xplore, ACM Digital Library (ACM DL) e Web of Science. O acesso é
pelo Portal de Periódicos CAPES, com a identificação institucional da UnB
(CAFe). No Web of Science, a busca usa só a \textit{Core Collection}. Na ACM
DL, a busca avançada, aberta pelo Portal CAPES, pesquisa a coleção
\textit{ACM Guide to Computing Literature}.
% PENDENTE: justificar a escolha de cada base. Os protocolos usam a descrição
% que cada empresa faz da própria base (Elsevier, IEEE, ACM, Clarivate). Essas
% páginas estão registradas em referencias-verificadas.md, mas não são
% referências revisadas por pares. Decidir se entram como nota de rodapé.

\textbf{Fontes nacionais.} As revisões também buscam em veículos brasileiros
que declaram revisão por pares na página oficial e que têm artigos publicados
sobre o tema da revisão. A Tabela~\ref{tab:fontes-nacionais} lista essas
fontes. Na RSL 1, a busca no SciELO é feita diretamente no site do SciELO.
Dois periódicos da Sociedade Brasileira de Computação, o \textit{Journal of
the Brazilian Computer Society} (JBCS) e o \textit{Journal of Internet Services
and Applications} (JISA), não têm busca separada, porque estão cobertos pelo
Scopus. Só os volumes do JBCS de 1997 a 1999 ficam fora do Scopus. Na RSL 2, os anais do SBES entram só a partir de 2024, porque os de
2017 a 2023 foram publicados pela ACM e já entram na busca da ACM DL. O JSERD
tem busca separada, porque não está no Scopus.

\begin{table}[htbp]
	\centering
	\caption{Fontes nacionais das revisões}
	\label{tab:fontes-nacionais}
	\small
	\begin{tabular}{p{1.1cm}>{\raggedright\arraybackslash}p{5.6cm}>{\raggedright\arraybackslash}p{2.8cm}>{\raggedright\arraybackslash}p{4.6cm}}
		\hline
		\textbf{RSL} & \textbf{Fonte} & \textbf{Editora} &
		\textbf{Revisão por pares declarada} \\
		\hline
		1 & Revista Brasileira de Inteligência (RBI) & ABIN &
		Duplo-cega, exceto para artigos depositados em servidor de
		\textit{preprints} \\
		1 & Anais do SBSeg (Simpósio Brasileiro de Cibersegurança) & SBC &
		Duplo-cega, em duas fases (revisão e meta-revisão) \\
		1 & Revista Brasileira de Segurança Pública (RBSP) & Fórum Brasileiro
		de Segurança Pública & Duplo-cega \\
		1 & Revista Brasileira de Direito Processual Penal (RBDPP) & IBRASPP &
		Duplo-cega \\
		1 & Revista Brasileira de Criminalística (RBC) & Associação Brasileira
		de Criminalística & Duplo-cega, com pelo menos dois pareceristas
		\textit{ad hoc} \\
		1 & SciELO (todas as coleções) & Não se aplica & Depende do periódico \\
		\hline
		2 & Anais do SBES (Simpósio Brasileiro de Engenharia de Software),
		de 2024 em diante & SBC & Duplo-anônima, com pelo menos três membros
		do comitê de programa por artigo \\
		2 & \textit{Journal of Software Engineering Research and Development}
		(JSERD) & SBC & \textit{Single-anonymous}, com pelo menos dois
		revisores \\
		\hline
	\end{tabular}
	\fonte{Elaborada pelo autor, com base nos protocolos das revisões.}
\end{table}
% PENDENTE: as evidências de revisão por pares e de indexação de cada fonte
% nacional (páginas das revistas, DOAJ, SciELO, Qualis na Plataforma
% Sucupira) estão em referencias-verificadas.md e nos protocolos, mas não são
% referências bibliográficas. Decidir se entram como notas de rodapé ou como
% apêndice. A RBC entra pelo Qualis A2 (quadriênio 2021-2024), porque não
% está no DOAJ e não declara Scopus, Web of Science ou SciELO.

\textbf{Busca complementar.} As duas revisões fazem \textit{snowballing}
\textit{backward} e \textit{forward} a partir dos estudos incluídos, repetido
até que uma iteração não traga nenhum estudo novo \cite[seção 3.3 e
Figura~1]{wohlin2014}. O \textit{forward} usa o Google Scholar
\cite[seção 3.2.2]{wohlin2014}. Na RSL 2, os estudos incluídos nas três
revisões citadas acima também são conferidos e passam pela mesma triagem, e o
relatório explica como essas revisões contribuíram para a busca, como pede o
SEGRESS \cite[item 7]{kitchenham2023segress}.

% ---------------------------------------------------------------------------
\subsection{Estratégia de busca}

\textbf{RSL 1.} A string combina dois blocos de termos. O bloco A reúne os
termos de OSINT:

{\footnotesize
\begin{verbatim}
"open source intelligence" OR "open-source intelligence" OR OSINT OR SOCMINT
OR "social media intelligence"
\end{verbatim}
}

\noindent O bloco B reúne os termos do contexto policial e de inteligência:

{\footnotesize
\begin{verbatim}
police OR policing OR "law enforcement" OR "criminal investigation*"
OR "digital investigation*" OR "cyber investigation*" OR "intelligence agenc*"
OR "intelligence analys*" OR "public security" OR crime OR forensic*
OR "national security" OR "intelligence communit*" OR "intelligence service*"
OR "security communit*" OR "national intelligence" OR counterintelligence
OR "counter-intelligence"
\end{verbatim}
}

\noindent A string busca \texttt{A AND B} em título, resumo e palavras-chave,
ou o bloco A só no título. No Scopus, ela foi executada assim:

{\footnotesize
\begin{verbatim}
TITLE-ABS-KEY ( ( "open source intelligence" OR "open-source intelligence"
  OR osint OR socmint OR "social media intelligence" )
AND ( police OR policing OR "law enforcement" OR "criminal investigation*"
  OR "digital investigation*" OR "cyber investigation*"
  OR "intelligence agenc*" OR "intelligence analys*" OR "public security"
  OR crime OR forensic* OR "national security" OR "intelligence communit*"
  OR "intelligence service*" OR "security communit*"
  OR "national intelligence" OR counterintelligence
  OR "counter-intelligence" ) )
OR TITLE ( "open source intelligence" OR "open-source intelligence"
  OR osint OR socmint OR "social media intelligence" )
\end{verbatim}
}

% PENDENTE: o PRISMA-S (item 8) pede a estratégia de cada base exatamente como
% foi executada. Incluir em apêndice as strings do Web of Science, as 4
% sub-buscas do IEEE Xplore e a string da ACM DL (docs/rsl1-protocolo.md,
% seção 7.2; docs/buscas/ieee-xplore-2026-11-09.md;
% docs/buscas/acm-dl-2026-11-09.md). A string do Scopus acima foi quebrada em
% linhas para caber na página; conferir se o autor prefere a forma original.

O piloto, em 09/11/2026, mudou a string. A versão anterior, só com
\texttt{A AND B} e sem os termos de segurança nacional e de comunidade de
inteligência, não recuperava dois dos quatro estudos primários de controle:
um deles usa termos que não estavam no bloco B, e o outro não tem resumo no
Scopus. O bloco B foi ampliado, e entrou o ramo com o bloco A no título.

\textbf{RSL 2.} A string combina os blocos abaixo. O bloco A1 reúne os termos
de assistentes e agentes de código:

{\footnotesize
\begin{verbatim}
"coding assistant*" OR "programming assistant*" OR "AI pair programm*"
OR "AI-assisted programming" OR "AI-assisted software development"
OR "AI-assisted coding" OR "coding agent*" OR "software engineering agent*"
OR "LLM-based agent*" OR "LLM-based assistant*" OR "LLM assistant*"
OR Copilot OR "AI coding"
\end{verbatim}
}

\noindent O bloco A2 reúne os termos gerais de LLMs, e o bloco B, os de
desenvolvimento:

{\footnotesize
\begin{verbatim}
A2: "large language model*" OR LLM OR LLMs OR ChatGPT OR "generative AI"
B:  "software development" OR "software engineering" OR programming
    OR programmer* OR developer* OR "code generation" OR "generated code"
\end{verbatim}
}

\noindent O bloco C reúne os termos das dimensões estudadas:

{\footnotesize
\begin{verbatim}
produtividade: productiv* OR efficien*
qualidade:     "code quality" OR "software quality" OR defect* OR bug*
               OR maintainab*
segurança:     secur* OR vulnerab*
custo:         cost* OR token* OR economic* OR energy
processo:      "software process" OR "development process" OR workflow*
               OR governance OR "human oversight" OR "code review"
\end{verbatim}
}

\noindent Dois blocos só se aplicam ao título. O bloco T reúne os termos de
desenvolvedor, e o bloco K, os de código:

{\footnotesize
\begin{verbatim}
T: developer* OR programmer* OR "software development" OR "software engineer*"
K: code OR software
\end{verbatim}
}

\noindent A string tem três ramos, ligados por \texttt{OR}:
\texttt{A1 AND B} em título, resumo e palavras-chave; \texttt{A2 AND C} em
título, resumo e palavras-chave, com o bloco T no título; e o ramo de custo,
com \texttt{A2} e os termos de custo em título, resumo e palavras-chave e o
bloco K no título. A busca fica restrita a artigos de periódico e de evento.

% PENDENTE: incluir em apêndice as strings da RSL 2 por base, como foram
% executadas (docs/rsl2-protocolo.md, seção 8.2; sub-buscas do IEEE Xplore em
% docs/buscas/ieee-xplore-2026-11-09.md e ieee-xplore-2026-11-10.md; ACM DL em
% docs/buscas/acm-dl-2026-11-09.md e acm-dl-2026-11-10.md), depois das buscas
% definitivas.

O piloto também mudou a string da RSL 2. A versão anterior retornou 11.236
registros no Scopus. Restringir a busca aos termos de assistente e agente
(A1), mantendo o bloco C, reduziu o volume, mas perdeu dois estudos de
controle, que tratam de usabilidade e de interação, e não de produtividade ou
qualidade. Por isso o ramo A1 dispensa o bloco C. O ramo de custo entrou
porque a string anterior não recuperava o estudo de controle sobre custo, cujo
título fala de código, e não de desenvolvedores. Entre as duas revisões da
string medidas no IEEE Xplore, o autor escolheu o ramo de custo, que acrescenta
menos registros que ampliar o bloco T.

\textbf{Adaptações por base.} No IEEE Xplore, cada termo é buscado nos campos
de título, resumo e termos de indexação, porque, no piloto, aplicar um campo a
um grupo de termos retornou zero resultados. Com isso, a string passa do
limite de 10 curingas por busca do
IEEE Xplore\footnote{\url{https://ieeexplore.ieee.org/Xplorehelp/searching-ieee-xplore/search-tips}}.
Ela foi dividida em sub-buscas de no máximo 9 curingas (4 na RSL 1 e 21 na
RSL 2), e o resultado é a união delas, sem duplicados. Na ACM DL, o curinga
não funciona dentro de frase entre aspas, o que o piloto confirmou. Por isso
as frases com curinga são escritas com as variantes por extenso. Na RSL 2, os
três ramos foram executados separadamente na ACM DL, porque a string inteira
voltou com a página de resultados incompleta, e as listas foram unidas pelo
DOI.
% PENDENTE: a limitação do curinga em frases da ACM DL vem das dicas de busca
% da própria ACM DL (referencias-verificadas.md, piloto da ACM DL), que não
% são referência revisada por pares. Decidir se entra como nota de rodapé.

\textbf{Fontes nacionais.} Na RSL 1, a busca usa a caixa de busca de cada
site, com os cinco termos do bloco A e dois termos em português,
\lq\lq fontes abertas\rq\rq{} e \lq\lq inteligência de fontes abertas\rq\rq,
executados separadamente. As duas grafias de \lq\lq open source
intelligence\rq\rq{} ficam separadas, porque a busca dos sites não as trata
como iguais. Na RSL 2, a busca usa os termos dos blocos A1 e A2. Em todas as
fontes, a string é adaptada ao site e registrada exatamente como foi
executada, com a data \cite[itens 8 e 13]{rethlefsen2021}.

\textbf{Filtros, período e idioma.} As revisões não aplicam filtro de idioma
nas bases. O idioma é conferido na triagem (critério CI4), e a ausência de
limite é informada, como pede o PRISMA-S \cite[item 9]{rethlefsen2021}. Não há
limite inferior de data: a busca vai até a data da busca definitiva. Qualquer
restrição de data, idioma ou fonte é registrada com a justificativa
\cite[seção 4.3.1]{kitchenham2023segress}. A RSL 1 não aplica filtro de tipo
de documento. A RSL 2 aplica o filtro de artigos de periódico e de evento.

\textbf{Validação da string.} O PRISMA 2020 cita, como forma de validar a
busca, verificar se ela identifica um conjunto de estudos claramente elegíveis,
que deve ser informado \cite[item 7]{page2021ee}. As duas revisões usam
estudos de controle, encontrados por busca na Crossref, independente da
string. A RSL 1 tem quatro estudos primários de controle e confere também três
estudos secundários sobre OSINT
\cite{ghioni2024,browne2024,pastorgalindo2020}. A RSL 2 tem sete estudos
primários de controle, que cobrem as dimensões de segurança, qualidade,
processo, produtividade e custo. O autor (R1) e o orientador (R2) confirmam
na calibração que os controles atendem aos critérios de inclusão. Se algum
controle não for recuperado, a string é revisada, e a mudança é registrada no
histórico de versões do protocolo.
% PENDENTE: decidir se os estudos de controle (4 da RSL 1 e 7 da RSL 2) são
% listados em apêndice. Os metadados deles foram conferidos só na Crossref, e
% a revisão por pares de cada veículo ainda não foi conferida
% (referencias-verificadas.md, "Controles: metadados conferidos na Crossref").
% Por isso eles não foram citados neste capítulo.

A Tabela~\ref{tab:piloto} mostra o resultado do piloto com as strings atuais
de cada revisão.

\begin{table}[htbp]
	\centering
	\caption{Piloto das buscas com as strings atuais}
	\label{tab:piloto}
	\small
	\begin{tabular}{p{2.6cm}p{1.9cm}>{\raggedright\arraybackslash}p{3.4cm}p{1.9cm}>{\raggedright\arraybackslash}p{3.4cm}}
		\hline
		& \multicolumn{2}{l}{\textbf{RSL 1 (09/11/2026)}} &
		\multicolumn{2}{l}{\textbf{RSL 2 (10/11/2026)}} \\
		\textbf{Base} & \textbf{Registros} & \textbf{Controles recuperados} &
		\textbf{Registros} & \textbf{Controles recuperados} \\
		\hline
		Scopus & 882 & 4 de 4 primários e 3 de 3 secundários & 3.148 &
		6 de 6 indexados \\
		Web of Science & 494 & 4 de 4 primários e 3 de 3 secundários &
		1.909 & 5 de 5 indexados \\
		IEEE Xplore & 267 & 1 de 1 indexado & 1.577 & 2 de 2 indexados \\
		ACM DL & 137 & 2 de 2 indexados & 1.413 & 6 de 6 indexados \\
		\hline
	\end{tabular}
	\fonte{Elaborada pelo autor, com base nos protocolos das revisões.}
\end{table}

\textbf{Buscas definitivas da RSL 1.} As buscas definitivas da RSL 1 foram
feitas em 10/11/2026, com as strings do protocolo e sem filtro de data, de
idioma ou de tipo. As quatro bases retornaram o mesmo total do piloto:
882 registros no Scopus, 494 no Web of Science, 267 no IEEE Xplore e 137 na
ACM DL, num total de 1.780. Foram exportados 1.779 registros. O registro do
IEEE Xplore que não foi exportado é um curso, que a exportação de citações não
devolve. Nas fontes nacionais, a soma dos registros únicos de cada fonte deu 17, e
ficaram 15 depois de remover os duplicados entre fontes. Assim, a remoção de
duplicatas partiu de 1.794 registros.

\textbf{Teste de sensibilidade da string da RSL 1.} Como o protótipo se volta
para a investigação de corrupção e lavagem de dinheiro
(Seção~\ref{sec:metodologia-requisitos}), a string foi testada no Scopus, em
10/11/2026, com buscas que retornam só os registros que ela não recupera.
Termos de corrupção, lavagem de dinheiro e crimes financeiros, combinados com
o bloco A, trouxeram 8 registros fora da string. Termos como
\texttt{criminal*}, \texttt{terroris*} e \texttt{extremis*} trouxeram 60. A
união das duas listas tem 65 registros.
% PENDENTE: o teste de sensibilidade mostrou que o bloco B não tem
% "criminal" sozinho nem os termos de terrorismo e extremismo. Não há decisão
% registrada sobre ampliar a string ou tratar esses 65 registros (por exemplo,
% como busca complementar). Registrar a decisão no protocolo antes de
% descrevê-la aqui.

% ---------------------------------------------------------------------------
\subsection{Critérios de seleção}

Um estudo entra na revisão se atender a todos os critérios de inclusão (CI) e
a nenhum critério de exclusão (CE). As Tabelas~\ref{tab:criterios-rsl1}
e~\ref{tab:criterios-rsl2} apresentam os critérios de cada revisão.

\begin{table}[htbp]
	\centering
	\caption{Critérios de seleção da RSL 1}
	\label{tab:criterios-rsl1}
	\small
	\begin{tabular}{p{1cm}p{13.4cm}}
		\hline
		\textbf{ID} & \textbf{Critério} \\
		\hline
		\multicolumn{2}{l}{\textit{Inclusão (todos devem ser atendidos)}} \\
		CI1 & Trata de OSINT aplicada a investigação criminal, segurança
		pública ou atividade de inteligência \\
		CI2 & Apresenta ferramenta, sistema, arquitetura ou método, ou analisa
		requisitos, desafios ou riscos dessas ferramentas \\
		CI3 & Publicado em periódico ou evento com revisão por pares \\
		CI4 & Escrito em inglês ou português \\
		\hline
		\multicolumn{2}{l}{\textit{Exclusão (basta um)}} \\
		CE1 & OSINT apenas citada de passagem, sem ser objeto do estudo \\
		CE2 & OSINT exclusivamente comercial ou de marketing, sem finalidade
		investigativa ou de inteligência \\
		CE3 & Texto sem revisão por pares: \textit{preprint}, tese,
		dissertação, relatório técnico, editorial, resumo, pôster ou tutorial
		\\
		CE4 & Texto completo indisponível pela UnB/CAPES, por acesso aberto ou
		por pedido aos autores (pedido feito antes de excluir, com prazo
		registrado) \\
		CE5 & Duplicata (registro do mesmo relato) ou outro relato de um
		estudo já incluído \\
		\hline
	\end{tabular}
	\fonte{Elaborada pelo autor, com base no protocolo da RSL 1.}
\end{table}

\begin{table}[htbp]
	\centering
	\caption{Critérios de seleção da RSL 2}
	\label{tab:criterios-rsl2}
	\small
	\begin{tabular}{p{1cm}p{13.4cm}}
		\hline
		\textbf{ID} & \textbf{Critério} \\
		\hline
		\multicolumn{2}{l}{\textit{Inclusão (todos devem ser atendidos)}} \\
		CI1 & Estuda o uso de assistentes ou agentes baseados em LLMs em
		atividades de desenvolvimento de software \\
		CI2 & Traz evidência empírica (experimento, estudo de caso, estudo de
		campo, \textit{survey}, mineração de repositórios ou avaliação com
		\textit{benchmark}) sobre pelo menos uma das dimensões das QP1 a QP5
		\\
		CI3 & Publicado em periódico ou evento com revisão por pares \\
		CI4 & Escrito em inglês ou português \\
		\hline
		\multicolumn{2}{l}{\textit{Exclusão (basta um)}} \\
		CE1 & O LLM só é citado de passagem, sem ser objeto do estudo \\
		CE2 & A contribuição é só uma técnica para melhorar o desempenho do
		modelo em \textit{benchmark} (arquitetura, treino,
		\textit{fine-tuning}), sem analisar nenhuma das dimensões no
		desenvolvimento \\
		CE3 & Trata de engenharia de sistemas de IA (a IA como produto
		desenvolvido), e não do uso de IA para desenvolver software \\
		CE4 & Trata só de ensino, com foco em resultados de aprendizagem, ou
		todos os participantes humanos do estudo são estudantes. Estudos com
		amostra mista e estudos sem participantes humanos continuam elegíveis
		\\
		CE5 & Texto sem revisão por pares: \textit{preprint}, tese,
		dissertação, relatório técnico, editorial, resumo, pôster ou tutorial
		\\
		CE6 & Texto completo indisponível pela UnB/CAPES, por acesso aberto ou
		por pedido aos autores (pedido feito antes de excluir, com prazo
		registrado) \\
		CE7 & Duplicata (registro do mesmo relato) ou outro relato de um
		estudo já incluído \\
		\hline
	\end{tabular}
	\fonte{Elaborada pelo autor, com base no protocolo da RSL 2.}
\end{table}

Os relatos de um mesmo estudo são analisados juntos, com a regra de escolha
registrada \cite[quadro 1 e item 9]{page2021ee}. Nas fontes nacionais e em
eventos, o tipo do texto é conferido artigo por artigo na triagem, porque um
mesmo veículo pode ter mais de uma modalidade de revisão. É o caso da RBI, que
declara arbitragem duplo-cega, exceto para artigos depositados em servidor de
\textit{preprints}.

\textbf{Estudos secundários.} As revisões que atendam aos critérios não entram
como estudos primários. São analisadas à parte, como trabalhos relacionados, e
o relatório explica como contribuíram para a busca
\cite[item 7]{kitchenham2023segress}.

\textbf{Literatura cinzenta.} \citeonline{kitchenham2023grey} argumentam que
as diretrizes de revisão sistemática nunca excluíram a literatura cinzenta por
princípio: a decisão cabe aos critérios de elegibilidade de cada revisão, e um
estudo sem garantia de disponibilidade pública a longo prazo não deve ser
usado como estudo primário. As duas revisões deste trabalho exigem revisão por
pares (CI3), por decisão do autor. Na área da RSL 2, \citeonline[seção
2.3]{hou2024} mantiveram na revisão deles os \textit{preprints} do arXiv de
2023 e 2024, porque a área é emergente e muitos trabalhos ainda estavam em
submissão. O efeito dessa decisão sobre a cobertura é tratado nas ameaças à
validade (Seção~\ref{sec:metodologia-ameacas}).
% PENDENTE: Kitchenham, Madeyski e Budgen (2023b) foi conferido pela página do
% IEEE Xplore e pelo resumo (referencias-verificadas.md). Confirmar se o
% trecho sobre disponibilidade pública a longo prazo foi lido no texto
% completo.

\textbf{Regras de aplicação dos critérios da RSL 1.} A calibração de
10/11/2026 (Seção~\ref{sec:metodologia-selecao}) gerou regras para casos
ambíguos, aprovadas pelo autor (R1) e incorporadas ao protocolo:

\begin{itemize}
	\item livros e capítulos de livro não atendem ao CI3, que exige periódico
	ou evento; os relevantes podem ser usados na fundamentação do trabalho,
	fora da revisão;
	\item vigilância em saúde pública sem uso policial ou de inteligência,
	desastres e defesa civil não atendem ao CI1;
	\item cibersegurança corporativa, testes de invasão e ataques com OSINT
	contra empresas ou pessoas não atendem ao CI1, mas entram os estudos sobre
	ameaças contra a ferramenta ou contra o investigador (QP4);
	\item investigações de direitos humanos e jornalismo investigativo com
	fontes abertas atendem ao CI1 só quando a finalidade declarada for prova
	para responsabilização criminal;
	\item estudos que só aplicam OSINT a um alvo, sem propor ou avaliar
	ferramenta, método, requisito ou risco, não atendem ao CI2;
	\item registros sem resumo são decididos pelo título e pelo veículo e, se
	o título for compatível com CI1 e CI2, passam para a etapa 2;
	\item quando o idioma não pode ser identificado pelo registro, o registro
	passa para a etapa 2, e o CI4 é conferido no texto completo.
\end{itemize}

% PENDENTE: as regras foram aprovadas por R1 em 10/11/2026
% (docs/rsl1-calibracao.md, "Decisões"). Registrar se o orientador (R2)
% também as aprovou.

% ---------------------------------------------------------------------------
\subsection{Processo de seleção}
\label{sec:metodologia-selecao}

A seleção da RSL 1 tem sete passos:

\begin{enumerate}
	\item \textbf{Remoção de duplicatas}, antes da triagem, com registro de
	quantos registros saíram \cite[item 16a]{page2021ee}. A remoção é feita por
	um \textit{script} aplicado aos arquivos exportados das bases, fora do
	Rayyan, porque o recurso do Rayyan que resolve duplicatas em lote só está
	disponível nos planos
	pagos\footnote{\url{https://help.rayyan.ai/hc/en-us/articles/34563050458257}}.
	As regras são aplicadas em ordem: DOI igual; título normalizado e ano
	iguais, quando os DOIs não são diferentes; e conferência manual dos casos
	ambíguos. Entre os registros de um grupo de duplicatas, fica o que tem
	resumo e, entre esses, o da fonte nesta ordem: Scopus, Web of Science, IEEE
	Xplore, ACM DL e fontes nacionais.
	\item \textbf{Importação no Rayyan} \cite{ouzzani2016} do arquivo sem
	duplicatas. A detecção de duplicados do Rayyan serve de conferência, e
	cada par que ela marca é conferido à mão.
	\item \textbf{Calibração:} os dois revisores aplicam os critérios aos
	mesmos 20 a 30 registros, número recomendado por \citeonline[diretriz
	3]{polanin2019}, e discutem as divergências. Se necessário, os critérios
	são ajustados e registrados no histórico de versões do protocolo.
	\item \textbf{Etapa 1:} o autor (R1) e o orientador (R2) fazem a triagem
	de todos os registros por título, resumo e palavras-chave, de forma
	independente. Revisores independentes evitam o viés do pesquisador
	\cite[seção 2.2]{kitchenham2023segress}, e a triagem dupla tende a ser
	mais confiável que a simples \cite[item 8]{page2021ee}. O Rayyan fica no
	modo cego, que impede cada revisor de ver as decisões do outro até que os
	dois
	terminem\footnote{\url{https://help.rayyan.ai/hc/en-us/articles/17460918624145-How-to-Use-Blind-Mode-in-Rayyan}}.
	\item \textbf{Etapa 2:} o autor (R1) e o orientador (R2) leem o texto
	completo de todos os estudos
	que passaram na etapa 1, também de forma independente.
	\item \textbf{\textit{Snowballing}} iterativo, com a mesma triagem dupla e
	independente.
	\item \textbf{Registro do motivo de cada exclusão na etapa 2}, que o
	fluxograma do PRISMA 2020 apresenta para os relatos avaliados
	\cite{page2021prisma}.
\end{enumerate}

% PENDENTE: as duas notas de rodapé acima remetem à documentação do Rayyan.
% Em referencias-verificadas.md, a página do modo cego tem situação "⚠️"
% (documentação técnica, não revisada por pares), e a do Auto Resolver está
% registrada no registro de buscas. Decidir se a monografia mantém essas
% notas de rodapé.

\textbf{Concordância.} O SEGRESS pede que se informe como a concordância entre
os revisores foi medida e que ela seja relatada \cite[itens 8 e
16a]{kitchenham2023segress}. Em cada etapa, a medida é o kappa de Cohen
\cite{cohen1960}, calculado sobre todas as decisões. As divergências são
resolvidas por consenso entre o autor (R1) e o orientador (R2), e cada uma é
registrada.
% PENDENTE: se o capítulo de resultados interpretar o valor do kappa (por
% exemplo, por faixas de concordância), é preciso uma referência verificada
% para essas faixas.

\textbf{Diferenças na RSL 2.} A RSL 2 segue o mesmo processo. Pelo volume de
registros, ela usa a semiautomação do Rayyan, que ordena os registros com um
modelo de predição \cite{ouzzani2016}, seguindo a recomendação de usar uma
ferramenta de mineração de texto na triagem \cite[diretriz 6]{polanin2019}. A
ferramenta só ordena os registros: os dois revisores continuam decidindo todos
eles, como pede a diretriz 7 de \citeonline{polanin2019}.
A remoção de duplicatas da RSL 2 usa o mesmo \textit{script} e as mesmas
regras da RSL 1, e a detecção de duplicados do Rayyan serve de conferência.

% Execução da seleção (usado na monografia e nos PDFs das RSLs).
\subsubsection{Execução da seleção}

\textbf{Duplicatas e calibração.} Na RSL 1, a busca nas bases e a ampliação
do bloco B somaram 1.928 registros. Saíram 789 duplicatas, e ficaram 1.139
para a triagem. Na RSL 2, dos 8.181 registros saíram 3.858 duplicatas, e
ficaram 4.323. Antes da triagem dos revisores, os critérios foram aplicados a
25 registros sorteados de cada revisão, com semente fixa, numa aplicação
piloto feita com assistência de IA só a partir do registro, para testar se os
critérios estavam claros. Dela saíram as regras de aplicação apresentadas
acima. O autor (R1) e o orientador (R2) triaram os mesmos 25 registros de
forma independente. As decisões dos dois coincidiram nas duas amostras
(kappa = 1,0).
% PENDENTE: a soma dos arquivos da RSL 1 dá 1.928 registros, mas o registro
% de buscas diz "1.923" na seção de ampliação (docs/triagem/prisma-etapa2.md,
% "Observações"). Corrigir o registro de buscas antes de desenhar o
% fluxograma.

\textbf{Etapa 1.} R1 e R2 triaram todos os registros por título, resumo e
palavras-chave, no Rayyan, de forma independente e no modo cego. As decisões
dos dois coincidiram em todos os registros: 1.139 de 1.139 na RSL 1 e 4.323
de 4.323 na RSL 2 (kappa = 1,0 nas duas). Não houve divergência a resolver.
Passaram para a etapa 2 os registros incluídos ou em dúvida: 450 na RSL 1 e
1.115 na RSL 2.
% PENDENTE: concordância perfeita em 5.462 decisões é incomum. Guardar a
% exportação das decisões de cada revisor no Rayyan, para mostrar à banca
% que a triagem foi independente.

\textbf{Recorte da etapa 2.} O volume da etapa 2 é incompatível com a
leitura dupla e independente do texto completo. Por decisão do autor, entraram
critérios de inclusão aplicados só a partir da etapa 2. Mudanças nos
critérios depois do início da revisão são permitidas desde que relatadas com
o motivo \cite[item 24c]{page2021prisma}.
\begin{itemize}
	\item \textbf{RSL 1 (protocolo v0.16):} CI5, a ferramenta ou o sistema
	relata avaliação; CI6, o uso por órgão policial, de investigação criminal
	ou de inteligência é explícito no título, no resumo ou nas palavras-chave;
	CI7, publicado a partir de 2015.
	\item \textbf{RSL 2 (protocolo v0.12):} CI5, a evidência vem de
	desenvolvedores profissionais ou de amostra mista de profissionais e
	estudantes, e não só de \textit{benchmark}; CI6, periódico ou conferência
	da trilha principal com nota A* ou A no CORE 2023.
\end{itemize}
Os registros foram classificados pelo resumo e pelos metadados. Os casos
incertos foram para o texto completo, e os excluídos por esses critérios são
relatados numa caixa própria do fluxograma.

\textbf{Leitura do texto completo.} Os textos completos foram obtidos pelo
acesso institucional da UnB (Portal de Periódicos da CAPES) e em acesso
aberto. Os relatos que não foi possível obter são relatados como não
recuperados. A leitura do texto completo foi feita pela IA, como terceira
opinião, com as regras de aplicação R1-A a R1-D (RSL 1, v0.18) e R2-A a R2-H
(RSL 2, v0.14), propostas a partir das dúvidas da leitura e aceitas por R1.
% PENDENTE: R2 ainda confirma as regras R1-A a R1-D e R2-A a R2-H. A leitura
% independente de R1 e R2 na etapa 2, e o kappa dessa etapa, ainda não foram
% feitos. Se R2 mudar alguma regra, os números do fluxograma são recalculados
% por script (docs/triagem/recorte-etapa2/).

\textbf{\textit{Snowballing}.} Os estudos incluídos serviram de sementes para
o \textit{snowballing} \textit{backward} e \textit{forward}, repetido até que
uma iteração não trouxesse estudo novo \cite[seção 3.3]{wohlin2014}. A RSL 1
parou na sexta iteração, e a RSL 2, na terceira. No lugar do Google Scholar,
as referências e citações vieram do OpenAlex e da Semantic Scholar (RSL 1
v0.19, RSL 2 v0.15). Os registros novos passaram pelos mesmos critérios, com
exclusões automáticas pelos metadados (ano, idioma, tipo e veículo) e as
mesmas regras do recorte. A triagem e a leitura do texto completo também
foram feitas pela IA como terceira opinião, com as regras aceitas por R1.
% PENDENTE: R2 ainda confirma as regras do snowballing (R1-E a R1-H e as
% equivalentes da RSL 2).

\textbf{Declaração do uso de IA.} O modelo de linguagem Claude (Anthropic)
foi usado como assistente em todas as etapas das revisões: na aplicação
piloto da calibração, como terceira opinião na etapa 1, e na leitura do texto
completo, no \textit{snowballing}, na avaliação de qualidade, na extração e na
síntese. Em cada etapa, as instruções dadas ao modelo, as decisões e um
trecho literal do texto que sustenta cada decisão estão guardados no
repositório. Os trechos e os números foram conferidos por \textit{script}
contra o texto completo. Nenhuma regra nova foi aplicada sem a aprovação do
autor, e as decisões do modelo não substituem a avaliação dos revisores
prevista nos protocolos.
% PENDENTE: confirmar com o autor esta declaração e se há norma da UnB ou da
% FCTE sobre o uso de IA em trabalhos de conclusão.


% ---------------------------------------------------------------------------
% Avaliação de qualidade (usado na monografia e nos PDFs das RSLs).
\subsection{Avaliação de qualidade}
\label{sec:metodologia-qualidade}

As duas revisões usam o instrumento de \citeonline[seção 3.5]{dyba2008}: 11
critérios baseados no \textit{Critical Appraisal Skills Programme} (CASP) e em
boas práticas de pesquisa empírica em engenharia de software, que cobrem
rigor, credibilidade e relevância. Cada critério é respondido com sim ou não,
como no original. A Tabela~\ref{tab:qualidade} mostra os critérios e como
cada revisão os usa.
% PENDENTE: em referencias-verificadas.md, a linha da tabela de 27/10/2026
% ainda marca Dybå e Dingsøyr (2008) como "⚠️ retirada dos protocolos". A
% seção "Checklist de qualidade (atualização de 09/11/2026)" registra a
% leitura do texto completo e a revisão por pares. Atualizar a situação para
% ✅ naquela tabela.

\begin{table}[htbp]
	\centering
	\caption{Critérios de avaliação de qualidade}
	\label{tab:qualidade}
	\small
	\begin{tabular}{p{1.8cm}p{8.8cm}p{1.4cm}p{1.4cm}}
		\hline
		\textbf{Critério} & \textbf{Pergunta} & \textbf{RSL 1} &
		\textbf{RSL 2} \\
		\hline
		1 & O estudo relata pesquisa empírica, e não apenas lições aprendidas
		baseadas em opinião de especialista? & AQ1 & Não usa \\
		2 & Os objetivos estão claramente relatados, com a justificativa do
		estudo? & AQ2 & AQ1 \\
		3 & O contexto em que a pesquisa foi feita está adequadamente
		descrito? & AQ3 & AQ2 \\
		4 & O desenho da pesquisa é adequado aos objetivos? & AQ4 & AQ3 \\
		5 & A amostra e a forma de identificá-la e recrutá-la estão
		adequadamente descritas? & AQ5 & AQ4 \\
		6 & Houve grupo de controle para comparar tratamentos? & AQ6 & AQ5 \\
		7 & Os métodos de coleta de dados são adequados e estão descritos? &
		AQ7 & AQ6 \\
		8 & Os métodos de análise dos dados estão adequadamente descritos? &
		AQ8 & AQ7 \\
		9 & A relação entre pesquisador e participantes foi considerada de
		forma adequada? & AQ9 & AQ8 \\
		10 & Os achados estão claramente declarados, com resultados críveis e
		conclusões justificadas? & AQ10 & AQ9 \\
		11 & O estudo tem valor para a pesquisa ou para a prática? & AQ11 &
		AQ10 \\
		Acréscimo & O estudo informa o modelo e a versão (ou a data) da
		ferramenta usada? & Não usa & AQ11 \\
		\hline
	\end{tabular}
	\fonte{Elaborada pelo autor, com base em \citeonline[seção 3.5]{dyba2008}
	e nos protocolos das revisões.}
\end{table}

As duas revisões adaptam o instrumento, e adaptações de um instrumento
existente devem ser informadas \cite[item 11]{page2021ee}:

\begin{itemize}
	\item \textbf{RSL 1:} no original, o critério 1 serve para excluir estudos
	que não relatam pesquisa empírica. Aqui ele é só um item da nota, porque o
	CI2 admite propostas de ferramenta sem avaliação empírica.
	\item \textbf{RSL 2:} o critério 1 não entra na nota, porque o CI2 já
	exige evidência empírica. Entra um item sobre o modelo e a versão da
	ferramenta (AQ11), porque o desempenho e o comportamento do GPT-3.5 e do
	GPT-4 podem variar muito ao longo do tempo \cite{chen2024}.
\end{itemize}

Na RSL 1, o autor (R1) e o orientador (R2) avaliam cada estudo de forma independente e resolvem as
divergências por consenso. A pontuação de cada estudo, de 0 a 11, serve para
ponderar a síntese, e não para excluir estudos. Ponderar a importância de cada
estudo na síntese é um dos usos da avaliação de qualidade relatados na
engenharia de software, e o uso pretendido precisa ser declarado
\cite[seção 2.3 e tabela 2]{kitchenham2023segress}.
% PENDENTE: o protocolo da RSL 2 (seção 12) não diz se R1 e R2 avaliam a
% qualidade de forma independente. Confirmar e registrar no protocolo.

\textbf{Regras de aplicação.} O instrumento original não trata estudos sem
participantes humanos nem comparações com \textit{baseline}. Por decisão do
autor, quatro regras foram registradas nos protocolos (RSL 1 v0.20, RSL 2
v0.16):
\begin{itemize}
	\item \textbf{Q-A, amostra sem pessoas} (critério 5): sem participantes
	humanos, a amostra é o conjunto de dados avaliado. A resposta é sim quando
	o texto descreve a origem dos dados e como foram selecionados.
	\item \textbf{Q-B, relação entre pesquisador e participantes} (critério 9):
	sim quando o texto discute como os pesquisadores podem ter influenciado os
	dados ou os resultados, por exemplo ao avaliar a própria ferramenta.
	Aprovação ética, consentimento ou declaração padrão de conflito de
	interesse, sozinhos, não bastam.
	\item \textbf{Q-C, grupo de controle} (critério 6): sim quando há uma
	condição de comparação medida no mesmo estudo, inclusive um
	\textit{baseline} ou método concorrente rodado nos mesmos dados.
	\item \textbf{Q-D, modelo e versão} (AQ11 da RSL 2): sim quando o estudo
	nomeia o modelo ou o produto e a versão, ou a data de uso.
\end{itemize}

\textbf{Execução.} A avaliação de qualidade e a extração foram feitas na
mesma leitura do texto completo, pela IA como terceira opinião, para não ler
cada texto duas vezes. Cada resposta traz um trecho literal do estudo e o
local, e um \textit{script} conferiu que todos os trechos estão no texto. Um
piloto com 8 estudos mostrou que a regra Q-B estava frouxa. Ela foi apertada
e os 8 estudos foram reavaliados.
% PENDENTE: a avaliação independente de R1 e R2 e a conferência da extração
% por R2 ainda não foram feitas. As regras Q-A a Q-D aguardam R2.


% ---------------------------------------------------------------------------
\subsection{Extração de dados}

Nas duas revisões, o autor (R1) extrai os dados de todos os estudos incluídos,
e o orientador (R2) confere a extração de todos eles, uma das formas de coleta
descritas no PRISMA 2020 \cite[item 9]{page2021ee}. Os formulários de extração
têm campos comuns (referência completa, ano, país, tipo de publicação, tipo de
estudo e nota de qualidade) e campos ligados às questões de pesquisa:

\begin{itemize}
	\item \textbf{RSL 1:} ferramenta ou sistema descrito e sua arquitetura e
	fontes de dados OSINT utilizadas (QP1); técnicas de IA e fase do ciclo de
	inteligência (QP2); requisitos e desafios legais, éticos, de privacidade e
	de soberania (QP3); ameaças, riscos e mitigações (QP4); método de
	avaliação, métricas e tratamento da evidência, isto é, preservação, cadeia
	de custódia e rastreabilidade (QP5). As fases do ciclo de inteligência são
	planejamento, coleta, processamento, análise, difusão e avaliação,
	segundo \citeonline[seção 1]{browne2024}.
	\item \textbf{RSL 2:} número de participantes ou projetos; contexto e
	tarefa de desenvolvimento, ferramenta, modelo, versão e nível de autonomia
	(QP6); grupo de comparação, se houver; métricas e resultados de
	produtividade (QP1) e de qualidade (QP2); vulnerabilidades, riscos e
	mitigações avaliadas (QP3); custos relatados e forma de medição (QP4); e
	mudanças de processo, papéis, revisão humana e governança (QP5).
\end{itemize}

% ---------------------------------------------------------------------------
\subsection{Síntese}

As duas revisões começam por um mapeamento, isto é, a análise das
características dos estudos \cite[seção 4.2]{kitchenham2023segress}. Na RSL 1,
o mapeamento mostra a distribuição dos estudos por ano, país, tipo de estudo,
técnica de IA e fase do ciclo de inteligência. Na RSL 2, por ano, contexto,
tipo de estudo, ferramenta e nível de autonomia.

Na RSL 1, cada questão de pesquisa recebe uma síntese temática, nas três
etapas de \citeonline{thomas2008}: codificação do texto linha a linha,
desenvolvimento de temas descritivos e geração de temas analíticos.

Na RSL 2, cada dimensão (QP1 a QP5) recebe uma tabela com a direção do efeito
relatado (positivo, negativo, misto ou nulo), as métricas usadas e a nota de
qualidade de cada estudo. A contagem pela direção do efeito é uma das formas
de síntese descritas no PRISMA 2020 \cite[quadro 1]{page2021ee}. Não está
prevista meta-análise: \citeonline{mohamed2026} descrevem os estudos sobre
assistentes baseados em LLMs como metodologicamente diversos, em sua maioria
exploratórios (59\%), e com resultados contraditórios sobre a qualidade, que
dependem do contexto e do critério de avaliação. Se o corpus permitir, essa
decisão é revista e registrada.
% PENDENTE: o trecho de Mohamed, Assi e Guizani (2026) foi conferido só no
% resumo, porque o texto completo não está em HTML na ACM DL
% (referencias-verificadas.md, "Sem meta-análise na RSL 2 (resumo)").

Nas duas revisões, a síntese termina na rastreabilidade para o trabalho: cada
achado relevante da RSL 1 vira um requisito candidato do buscador, e cada
achado relevante da RSL 2 vira uma prática ou um requisito candidato para o
processo de desenvolvimento.

% Regras de aplicação da síntese (usado na monografia e nos PDFs das RSLs).
\textbf{Regras de aplicação da síntese.} Por decisão do autor, os protocolos
registram regras para a síntese (RSL 1 v0.21, RSL 2 v0.17):
\begin{itemize}
	\item \textbf{O que é codificado:} a codificação linha a linha é feita
	sobre os campos da extração, que resumem o texto completo, e não sobre o
	próprio texto, como em \citeonline{thomas2008}. Os números dos campos de
	resultado foram conferidos por \textit{script} contra o texto completo.
	\item \textbf{Unidade de análise:} o estudo, e não o relato
	\cite[quadro 1]{page2021ee}. Relatos do mesmo estudo contam uma vez.
	\item \textbf{Peso da qualidade:} as contagens usam todos os estudos e
	são repetidas só com os de nota maior ou igual a 8, como análise de
	sensibilidade.
	\item \textbf{Efeito medido e percebido (RSL 2):} a direção do efeito
	sobre produtividade e qualidade também é contada pelo tipo de medida
	(métrica objetiva, percepção ou ambas), classificado estudo a estudo.
	\item \textbf{Síntese temática também na RSL 2:} a QP5 não tem direção de
	efeito. Por isso as questões da RSL 2 também passam pelas três etapas de
	\citeonline{thomas2008}.
	\item \textbf{Categorias do mapeamento:} país, tipo de estudo, técnica de
	IA, contexto e nível de autonomia vêm de texto livre e são normalizados
	por dicionários publicados com o \textit{script} do mapeamento.
\end{itemize}
A síntese também foi feita pela IA como terceira opinião. Cada código lista
os estudos que o apoiam, e um \textit{script} conferiu que os estudos citados
existem e têm informação no campo correspondente.
% PENDENTE: R2 ainda confirma as regras da síntese.


% ---------------------------------------------------------------------------
\subsection{Ameaças à validade}
\label{sec:metodologia-ameacas}

\begin{itemize}
	\item \textbf{Viés de publicação.} \citeonline[p.~10302]{pastorgalindo2020}
	descrevem o uso de OSINT pelos órgãos policiais e de inteligência da
	Espanha como um
	aspecto confidencial do funcionamento desses órgãos, o que pode deixar
	ferramentas sem publicação. Exigir revisão por pares também deixa de fora
	o que só é publicado fora desses veículos, como os \textit{preprints} que
	\citeonline[seção 2.3]{hou2024} mantiveram na área da RSL 2. Buscas amplas
	são a defesa contra o viés de publicação
	\cite[seção 2.2]{kitchenham2023segress}. Na RSL 1, a mitigação prevista é
	triangular os resultados com a entrevista com um delegado de polícia.
	\item \textbf{Terminologia.} Termos inconsistentes fazem a busca em bases
	perder estudos relevantes, como no exemplo de \citeonline[seção
	3.1]{wohlin2014}, e termos amplos trazem muitos artigos irrelevantes
	\cite[seção 2]{wohlin2014}. O piloto calibra a string, e os estudos de
	controle verificam se ela recupera o que deve.
	\item \textbf{Subjetividade da seleção.} É mitigada pela triagem dupla e
	independente \cite[seção 2.2]{kitchenham2023segress}, pela calibração e
	pela medida de concordância.
	\item \textbf{Idioma e acesso.} As revisões só incluem textos em inglês e
	português, e produção relevante em outros idiomas pode ficar de fora. A
	restrição de idioma e a impossibilidade de acessar todos os relatos são
	limitações a discutir \cite[item 23c]{page2021ee}.
	\item \textbf{Obsolescência (RSL 2).} Um resultado obtido com uma versão do
	modelo pode não valer para a seguinte \cite{chen2024}. O item AQ11 e a
	extração da versão do modelo permitem relativizar os achados.
\end{itemize}

% ===========================================================================
\section{Registro e transparência}
\label{sec:metodologia-registro}

Os protocolos das duas revisões não foram registrados em base externa. Eles
ficam versionados no repositório do trabalho, cada um com um histórico de
versões que registra as mudanças, com data e motivo. Na data descrita neste
capítulo, o protocolo da RSL 1 está na versão 0.22, e o da RSL 2, na versão
0.18. O PRISMA 2020 pede que o relato informe o registro da revisão ou declare
que ela não foi registrada, e que indique onde o protocolo pode ser acessado
\cite[itens 24a e 24b]{page2021prisma}. No SEGRESS, o registro é opcional para
todos os tipos de revisão \cite[item 24a]{kitchenham2023segress}. O documento
de explicação e elaboração do PRISMA 2020 associa o registro à redução de viés
e ao aumento da transparência \cite[item 24a, quadro 6]{page2021ee}.
% PENDENTE: indicar o endereço público do repositório (ou anexar os
% protocolos como apêndice), para atender ao item 24b do PRISMA 2020.

As buscas seguem o PRISMA-S. Para cada base, o registro de buscas guarda o
nome e a plataforma, a data, a string exata executada, os filtros aplicados e
o número de resultados \cite[itens 1, 8, 9, 13 e 15]{rethlefsen2021}. Os
arquivos exportados das bases ficam guardados localmente, fora do controle de
versões. A remoção de duplicatas gera um relatório com cada grupo de
duplicatas, o registro mantido e os removidos. O sorteio da amostra de
calibração usa semente fixa e pode ser repetido.
% PENDENTE: os arquivos exportados (docs/exportacoes/) ficam fora do git.
% Decidir se e onde serão disponibilizados (por exemplo, em repositório de
% dados), já que o PRISMA-S trata do relato das buscas, e não da guarda dos
% resultados.

Cada referência citada nos protocolos foi conferida na página do editor, e
cada escolha de método cita a fonte e o local do trecho que a sustenta, como
nas citações deste capítulo.

% ===========================================================================
\section{Levantamento de requisitos}
\label{sec:metodologia-requisitos}

O levantamento de requisitos partiu de uma sugestão do orientador, recebida em
10/11/2026, para um aplicativo OSINT de investigação de corrupção e lavagem de
dinheiro na Polícia Civil do Distrito Federal (PCDF). Esse levantamento é uma
fonte primária do projeto, e não uma referência bibliográfica. Ele tem dois
eixos: correlacionar dados abertos heterogêneos (pessoas, empresas, contratos,
bens, redes) e produzir prova que se sustente em juízo.

A lista inicial tem 10 requisitos funcionais (RF01 a RF10) e 11 requisitos não
funcionais (RNF01 a RNF11). Cada requisito foi ligado à questão de pesquisa da
RSL 1 que pode trazer evidência sobre ele, como mostra a
Tabela~\ref{tab:requisitos}. Para o orientador, o requisito mais crítico para
a defesa em juízo é a combinação da cadeia de custódia (RF07) com a trilha de
auditoria (RNF02).

\begin{table}[htbp]
	\centering
	\caption{Lista inicial de requisitos e ligação com a RSL 1}
	\label{tab:requisitos}
	\small
	\begin{tabular}{p{1.4cm}p{9.6cm}p{2.6cm}}
		\hline
		\textbf{ID} & \textbf{Requisito} & \textbf{QP da RSL 1} \\
		\hline
		RF01 & Coleta multifonte & QP1 \\
		RF02 & Resolução de entidades & QP1, QP2 \\
		RF03 & Análise de vínculos em grafo & QP1, QP2 \\
		RF04 & Alertas de tipologias & QP2 \\
		RF05 & Linha do tempo patrimonial & QP1 \\
		RF06 & Monitoramento contínuo & QP1 \\
		RF07 & Cadeia de custódia digital & QP5 \\
		RF08 & Gestão de casos & QP1 \\
		RF09 & Relatórios & QP5 \\
		RF10 & Integrações sob ordem judicial & QP3 \\
		\hline
		RNF01 & Segurança & QP3, QP4 \\
		RNF02 & Auditoria imutável & QP4, QP5 \\
		RNF03 & Anonimato operacional & QP4 \\
		RNF04 & Conformidade legal & QP3 \\
		RNF05 & Explicabilidade & QP2, QP5 \\
		RNF06 & Desempenho e escala & QP5 \\
		RNF07 & Disponibilidade e resiliência & QP4 \\
		RNF08 & Usabilidade & QP5 \\
		RNF09 & Interoperabilidade & QP1 \\
		RNF10 & Manutenibilidade & QP1 \\
		RNF11 & DevSecOps & QP4 \\
		\hline
	\end{tabular}
	\fonte{Elaborada pelo autor, com base no levantamento com o orientador.}
\end{table}

O levantamento segue três passos:

\begin{enumerate}
	\item \textbf{Confronto com a evidência e com a legislação.} Cada requisito
	é confrontado com a evidência da RSL 1, pela questão de pesquisa indicada
	na Tabela~\ref{tab:requisitos}, e com a legislação conferida. O resultado
	desse confronto entra na monografia.
	\item \textbf{Conferência das fontes de dados.} Para cada fonte do RF01, o
	autor confere se há API ou dados abertos oficiais, os termos de uso e quais
	campos são públicos.
	\item \textbf{Recorte do protótipo.} O autor define quais requisitos entram
	na implementação e quais ficam como trabalho futuro. Os requisitos são
	escritos como histórias de usuário, agrupados em épicos e priorizados pelo
	método MoSCoW (Capítulo~\ref{cap:requisitos}).
\end{enumerate}

% PENDENTE: citar a técnica de levantamento de requisitos usada (por exemplo,
% entrevista com especialista). Não há referência verificada sobre engenharia
% de requisitos em docs/referencias-verificadas.md.
% PENDENTE: a entrevista com o delegado está prevista (docs/rsl1-protocolo.md,
% seção 14), mas a data, o roteiro e a necessidade de termo de consentimento
% ainda estão em aberto. Descrever o procedimento da entrevista quando for
% definido.
% PENDENTE: a legislação conferida (LGPD, Marco Civil da Internet, Leis
% 9.613/1998, 12.846/2013 e 12.850/2013) está em docs/requisitos-app.md e
% docs/normas-iso.md, mas não em docs/referencias-verificadas.md. Incluir as
% entradas BibTeX das leis antes de citá-las. O Marco Civil da Internet ainda
% não foi conferido (docs/requisitos-app.md, nota 2).

\endgroup
